\section{Desain Eksperimen}

Berdasarkan variabel penelitian, dataset, algoritma pembanding, dan prosedur pengukuran yang telah ditetapkan, dirancang empat eksperimen yang masing-masing dipetakan terhadap hipotesis penelitian sebagai operasionalisasi aktivitas \textit{Demonstration} dan \textit{Evaluation}.

Eksperimen 1 menguji H1 melalui dua faktor, yaitu tingkat dependensi antaratribut dan dimensi ruang solusi (D). 
Ketujuh algoritma dijalankan pada lima dataset ABAC Lab dalam kondisi stasioner pada tiga tingkat dependensi, yaitu rendah, sedang, dan tinggi. 
Metrik yang diukur meliputi \textit{F1-score}, WSC, dan \textit{coverage}. 
Dataset \textit{workforce} $(D=436)$ dan \textit{e-document} $(D=1.907)$ menjadi dataset utama karena memiliki keragaman pola yang memadai, sedangkan \textit{healthcare}, project-\textit{management}, dan \textit{university} $(D=46$–$54)$ digunakan sebagai konteks pelengkap karena tingkat duplikasinya lebih tinggi. 
H1 menduga manfaat pemodelan dependensi CoDA-EDA semakin terlihat ketika tingkat dependensi antaratribut meningkat.

Eksperimen 2a menguji H2 melalui pengukuran jejak memori pada eksekusi yang sama dengan Eksperimen 1. 
Metrik utama yang digunakan adalah $MemPeak_{alloc}$, sedangkan $MemPeak_{sys}$ digunakan sebagai konteks penggunaan memori pada tingkat sistem. 
Pengujian ini membandingkan kebutuhan memori absolut antaralgoritma, khususnya pada dataset berdimensi besar.

Eksperimen 2b melengkapi pengujian H2 dengan mengevaluasi bentuk penskalaan memori terhadap ukuran populasi atau populasi virtual (NP), yang divariasikan pada
$
NP\in{25,50,100,200,400}.
$
Representasi solusi, fungsi \textit{fitness}, dataset, dimensi (D), dan prosedur pengukuran dipertahankan identik. 
Setiap konfigurasi dijalankan dalam proses terpisah agar pengukuran puncak memori tidak dipengaruhi eksekusi sebelumnya.

Pengaruh $NP$ terhadap memori bergantung pada mekanisme algoritma. 
Pada BOA, ACO, dan UBK, peningkatan $NP$ menambah jumlah kandidat yang disimpan sehingga komponen memorinya secara konseptual tumbuh terhadap $O(NP\times D)$. 
Sebaliknya, pada CoDA-EDA dan keluarga \textit{compact}, $NP$ hanya mengatur granularitas pembaruan probabilitas sebesar $\pm1/NP$ tanpa mematerialisasi $NP$ individu. 
iBOA juga menggunakan statistik inkremental tanpa menyimpan populasi penuh. 
Oleh karena itu, sensitivitas memori terhadap $NP$ diringkas menggunakan

\begin{equation}
R_{\mathrm{NP}}
=
\frac{
\textit{MemPeak}_{\mathrm{alloc}}(NP=400)
}{
\textit{MemPeak}_{\mathrm{alloc}}(NP=25)
}.
\label{eq:rasio-memori-np}
\end{equation}
\label{first:mempeakalloc}

Nilai $R_{NP}$ mendekati 1 menunjukkan jejak memori relatif tidak sensitif terhadap $NP$, sedangkan nilai yang lebih besar menunjukkan pertumbuhan memori. 
H2 mengharapkan kurva CoDA-EDA, algoritma \textit{compact}, dan iBOA relatif datar, sedangkan BOA, ACO, dan UBK meningkat seiring $NP$. 
Dengan demikian, Eksperimen 2b memberikan bukti empiris terhadap perbedaan perilaku memori $O(D)$ dan $O(NP\times D)$.

Eksperimen 3 menguji H3 dengan membandingkan CoDA-EDA dan iBOA pada lima dataset ABAC Lab dalam skenario \textit{drift} rendah dan sedang. 
Waktu komputasi setiap siklus pembaruan dicatat secara terpisah untuk siklus yang hanya memperbarui parameter dan siklus yang juga memperbarui struktur. 
H3 menduga CoDA-EDA membutuhkan waktu pembaruan rata-rata lebih rendah karena pembelajaran struktur dilakukan secara bersyarat ketika \textit{drift} melampaui ambang.

Eksperimen 4 menguji H4 melalui dua strategi pembaruan CoDA-EDA, yaitu \textit{warm-start} dengan melanjutkan model sebelum \textit{drift} dan \textit{re-mining} penuh dengan menginisialisasi ulang proses setelah \textit{drift} terdeteksi. 
Waktu re-optimasi dan kualitas kebijakan pasca-\textit{drift} dibandingkan untuk kedua strategi. 
H4 menduga \textit{warm-start} dapat mencapai kualitas yang sebanding dengan \textit{re-mining} penuh dengan waktu pembaruan lebih rendah, terutama pada \textit{drift} berintensitas rendah.

Pada eksperimen utama, setiap kombinasi algoritma, dataset, dan skenario dijalankan menggunakan enam \textit{seed} pembangkitan dataset dan lima replikasi algoritma, sehingga menghasilkan 30 replikasi per kombinasi. 
Validasi eksternal menggunakan lima \textit{seed}. 
Replikasi tersebut digunakan untuk mengakomodasi variabilitas stokastik dan menyediakan data yang memadai bagi analisis statistik.

\begin{longtable}{|c|>{\raggedright\arraybackslash}p{2.2cm}|p{1.6cm}|>{\raggedright\arraybackslash}p{2.8cm}|>{\raggedright\arraybackslash}p{3cm}|>{\raggedright\arraybackslash}p{2.8cm}|}
\caption{Ringkasan Desain Eksperimen}
\label{tab:desain-eksperimen} \\
\hline
\textbf{No} & \textbf{Eksperimen} & \textbf{Hipotesis} & \textbf{Algoritma Dilibatkan} & \textbf{Dataset} & \textbf{Variabel Diukur} \\
\hline
\endfirsthead

\multicolumn{6}{c}{{\bfseries \tablename\ \thetable{} -- lanjutan dari halaman sebelumnya}} \\
\hline
\textbf{No} & \textbf{Eksperimen} & \textbf{Hipotesis} & \textbf{Algoritma Dilibatkan} & \textbf{Dataset} & \textbf{Variabel Diukur} \\
\hline
\endhead

\hline
\multicolumn{6}{r}{Lanjut ke halaman berikutnya} \\
\endfoot

\hline
\endlastfoot

% ==================== DATA BARIS ====================
1 & Kualitas kebijakan vs. korelasi dan dimensi & H1 & 7 kondisi (CoDA-EDA + 6 baseline) & Baris 1--5 (lima dataset; dua taraf D utama 436 dan 1.907 + tiga dataset kecil D=46--54 sebagai konteks), kondisi stasioner, 3 tingkat korelasi & Akurasi/F-score, WSC, coverage \\
\hline
2a & Jejak memori antar-algoritma & H2 & CoDA-EDA vs BOA, iBOA, ACO, UBK & Sama dengan Eksperimen 1 (reuse eksekusi) & MemPeak\_alloc (utama), MemPeak\_sys\label{first:mempeaksys} \\
\hline
2b & Kurva penskalaan memori & H2 & CoDA-EDA vs BOA, iBOA, ACO, UBK & Baris 1--2; $NP$ divariasikan (25--400) pada dua nilai D & MemPeak\_alloc sebagai fungsi $NP$ dan $D$ \\
\hline
3 & Efisiensi waktu pembaruan & H3 & CoDA-EDA vs iBOA & Baris 1--5, skenario drift rendah-sedang & Waktu komputasi per siklus pembaruan \\
\hline
4 & Kualitas \& efisiensi warm-start & H4 & CoDA-EDA (warm-start vs re-mining penuh) & Baris 1--5, skenario drift rendah-sedang & Waktu re-optimasi, kualitas kebijakan pasca-drift \\
\hline
5 & Validasi eksternal & H1-H4 & 7 kondisi (CoDA-EDA + 6 baseline) & Baris 6 (Smart Building IoT) & F-score relatif garis sepele, presisi, recall, coverage, memori, warm-start \\
\hline

\end{longtable}

Keterkaitan antara keempat hipotesis, eksperimen yang mengujinya, perbandingan yang dilakukan, metrik yang diukur, dan uji statistik yang digunakan divisualisasikan sebagai peta pada Gambar \ref{fig:peta_eksperimen}.

\begin{figure}[H]
    \centering
    \includegraphics[width=1\textwidth]{gambar_eksperimen.png}
    \caption{Peta Eksperimen: Pemetaan Hipotesis H1-H4 terhadap Perbandingan, Metrik, dan Uji Statistik}
    \label{fig:peta_eksperimen}
\end{figure}

Gambar \ref{fig:peta_eksperimen} menegaskan bahwa setiap hipotesis diuji oleh tepat satu eksperimen dengan perbandingan, metrik, dan uji statistik yang spesifik.
Eksperimen 1 dan 2 membandingkan tujuh kondisi algoritma sekaligus sehingga memakai uji omnibus Friedman dengan post-hoc Wilcoxon berkoreksi Holm-Bonferroni, sedangkan Eksperimen 3 dan 4 hanya membandingkan dua kondisi sehingga cukup memakai Wilcoxon signed-rank berpasangan langsung.
Peta ini juga memperlihatkan efisiensi rancangan pada Eksperimen 2 yang memakai ulang Eksperimen 1.